\section{Chiral charge on disks above the recognition threshold}
\label{sec:band-charge}

Here the physical metrics \(g_i\) are volume-one positive Yamabe
minimizers on one nonstandard smooth homotopy sphere. Write
\(e_i=E_+(g_i)=E_-(g_i)\to I<2T\). Let \(q_i\) denote the
volume-one modified Yamabe representative used in
Section~\ref{sec:endpoint-charge}, and let
\(\mu_i=dV_{q_i}\), \(\ell_i^2=\int L^2dV\). The measures
\(\sigma_i=L^2dV\) and \(d\lambda_i^\pm\) are conformally invariant.
In particular \(0<Y_i/6\le\ell_i\le\sqrt{4e_i/3}\).

\begin{theorem}[A partial chiral residual bound]\label{thm:band-charge}
Suppose \(\liminf_i(Y_i-6\ell_i/\sqrt{2})>0\). After passing to
a subsequence, for every \(\epsilon,\zeta,r_*>0\) there are finitely
many fixed radii \(r_j<r_*\) and, at all sufficiently large indices,
disjoint smoothly parametrized closed standard disks with union
\(U_i=\bigcup_j D_{ij}\). Each disk lies in a ball
\(B_{k_i}(x_{ij},r_j)\) of a positive-Ricci metric \(k_i\in[g_i]\)
on which the integral of \(|\Rm_{g_i}|_{\rm tens}^2dV_{g_i}\)
is less than \(\zeta\). On these same disks, for both chiralities,
\begin{equation}\label{eq:band-charge-residual}
 \lambda_i^\pm(U_i)\ge Y_i^2/24-e_i-\epsilon,\qquad
 \lambda_i^\pm(M\setminus U_i)\le 2e_i-Y_i^2/24+\epsilon
                                  \le 4(e_i-T)+\epsilon .
\end{equation}
In particular the hypothesis holds whenever \(I<4T/3\), and each
sign carries at least \(4T-3e_i-\epsilon\) on the disks. If
\(I<5T/4\), their raw complement has energy strictly below \(T\)
for both signs when the tolerance is sufficiently small.
\end{theorem}

\begin{proof}
The inequalities \(Y_i^2\ge48(2T-e_i)\) and
\(18\ell_i^2\le24e_i\) imply the asserted uniform margin for
\(I<4T/3\). Choose the smooth positive-Ricci metric in the class
given in Section~\ref{sec:physical-volume} and normalize its scale
so the least radius enclosing half of \(\sigma_i(M)\) is one.
Denote an attaining center by \(p_i\). The measure is absolutely
continuous on the smooth source, so the ball-mass continuity argument
there applies. Every unit ball and the exterior of \(B_1(p_i)\)
have \(\sigma_i\)-mass at most \(\ell_i^2/2\). For a cutoff
supported in either kind of set the ordinary Yamabe inequality in
\(q_i\), its scalar bound \(R_{q_i}=6L_{q_i}+c_i\le6L_{q_i}\),
and H\"older give
\[
 (Y_i-6\ell_i/\sqrt 2)\|f\|_{4,q_i}^2\le6\|df\|_{2,q_i}^2.
\]
The conformal fourth-power gradient bound of
\eqref{eq:physical-small-support} therefore applies with this floor
and with \(\mu_i\) on its right. Inward logarithmic cutoffs on
unit balls and outward cutoffs based at \(p_i\) give
\eqref{eq:physical-power} for \(\mu_i\). Choose a fixed \(D\) large
enough that \(\mu_i(B_D(p_i))>1/2\). Packing this ball by a bounded
number of quarter-balls finds one of uniformly positive
\(\mu_i\)-mass. Applying the displayed inequality on the
corresponding half-ball with a linear cutoff bounds the
\(k_i\)-volume of \(B_{D+1}(p_i)\) below. Relative volume comparison
now gives \eqref{eq:physical-noncollapse} on bounded windows.

Take a common pointed limit and local weak limits
\(\mu,\sigma,\nu\) for \(\mu_i,\sigma_i\) and
\(\nu_i=|\Rm_{g_i}|^2_{\rm tens}dV_{g_i}\), respectively.
The power estimates give tightness of \(\mu_i\) and
\(\mu(\{x\})=0\) everywhere. For \(q_i=w_i^2k_i\), positivity
of \(R_{k_i}\) gives
\(-\Delta_{k_i}w_i\le V_iw_i\) with \(V_i=L_{q_i}w_i^2\).
On every source ball \(H\), \(\|V_i\|_{2,k_i;H}^2=\sigma_i(H)\).
At each point with \(\sigma(\{x\})=0\) choose a fixed small
continuity radius in a bounded window. On the corresponding balls
the potential term in \eqref{eq:physical-moser} at \(p=4\) is
absorbed. The local \(L^8(k_i)\) bound on \(w_i\) follows from
\(\int w_i^4dV_{k_i}=1\). Thus \(\mu\) is absolutely continuous
relative to \(\mathcal H^4\) on a neighborhood of each such point
by the local \(L^2\) bound on the densities and noncollapsed volume
convergence. A countable subcover and the zero \(\mu\)-mass of the
countably many atoms of \(\sigma\) give \(\mu\ll\mathcal H^4\).
The regular points outside the atoms of \(\nu\) have full
\(\mu\)-measure. The finite disjoint ball selection and smooth
embeddings in Theorem~\ref{thm:volume-balls} apply there by
differentiation of this density. For every fixed \(d>0\) they
give disks with the specified radii and physical curvature ceilings,
and \(\mu_i(M\setminus U_i)<d\) for all large indices.

Finally the constant test in the ordinary Yamabe quotient gives
\(\int L_{q_i}d\mu_i\ge Y_i/6\). On the omitted set this integral
is at most \(\ell_i\sqrt d\). Since the total volume is one,
Cauchy--Schwarz on \(U_i\) and the pointwise spectral bound yield
\[
 (\lambda_i^++\lambda_i^-)(U_i)
 \ge \tfrac32\sigma_i(U_i)
 \ge \tfrac32(Y_i/6-\ell_i\sqrt d)_+^2 .
\]
Either opposite chirality has total mass \(e_i\). Subtract it and
choose \(d\) small using the uniform upper bounds for \(Y_i,\ell_i\).
This proves \eqref{eq:band-charge-residual}.
\end{proof}

These estimates bound the unmodified physical charge of the
complement. A fixed tolerance selects a fixed finite family size.
The strict path budget, the measured allowances of the actual
collars, and the positive sign of a replacement are still needed
when applying the comparison on this family.
